\section{Zadanie 13}
\autor{Oskar Prusik}
Niech \( Y \) będzie zmienną z rozkładu dwumianowego z parametrami \( n = 20 \) oraz \( p = \frac{1}{2} \).
Wyznacz \( P(8 \leq Y \leq 10) \) dokładnie.
Następnie rozpisz \( Y \) jako sumę indykatorów i wyznacz to samo prawdopodobieństwo
za pomocą Centralnego Twierdzenia Granicznego.
Różnica pomiędzy tymi dwoma wynikami wynika z faktu, iż przybliżamy rozkład dyskretny
z małym \( n \) za pomocą rozkładu ciągłego.
Sytuację można poprawić zauważając, że dla \( Y \) zachodzi
\( P(8 \leq Y \leq 10) = P(7.5 \leq Y \leq 10.5) \).
Wyznacz szukane prawdopodobieństwo po zastosowaniu tej poprawki oraz porównaj otrzymane wyniki numerycznie.

\subsection*{Rozwiązanie}

Niech \( Y \sim \mathrm{Bin}\left(n=20, p=\frac12\right) \). Wyznaczamy \( P(8 \le Y \le 10) \).

\subsection*{1. Obliczenie dokładne}
\[
	P(Y=k) = \binom{20}{k} \left(\frac12\right)^{20}, \quad k=0,\dots,20
\]
\[
	P(8 \le Y \le 10) = \sum_{k=8}^{10} \binom{20}{k} \cdot \frac{1}{2^{20}}
\]
\[
	\binom{20}{8}=125970,\ \binom{20}{9}=167960,\ \binom{20}{10}=184756
\]
\[
	P_{\text{exact}} = \frac{125970+167960+184756}{2^{20}} = \frac{478686}{1048576} \approx 0.4565
\]

\subsection*{2. Przybliżenie CTG bez poprawki ciągłości}

Niech \( X_1,\dots,X_{20} \) i.i.d., \( X_i \sim \mathrm{Bernoulli}(1/2) \).
Wtedy:
\[
	\ev{X_i} = \mu = \frac12,\quad
	\variance{X_i} = \sigma^2 = \frac14,\quad
	\sigma = \frac12.
\]

\[
	\bar{X}_{20} = \frac{1}{20}\sum_{i=1}^{20} X_i = \frac{Y}{20}
\]

Z Twierdzenia 15.3.1:
\[
	\lim_{n\to\infty} P\!\left( a \le \frac{\bar{X}_n-\mu}{\sigma/\sqrt{n}} \le b \right) = \Phi(b)-\Phi(a)
\]
Dla \( n=20 \) stosujemy przybliżenie.

Zdarzenie \( \{8 \le Y \le 10\} \) równoważne jest \( \{0.4 \le \bar{X}_{20} \le 0.5\} \).

Podstawiamy do wzoru CTG:
\[
	a = \frac{0.4 - \mu}{\sigma/\sqrt{20}} = \frac{0.4 - 0.5}{0.5/\sqrt{20}}
\]
\[
	b = \frac{0.5 - \mu}{\sigma/\sqrt{20}} = \frac{0.5 - 0.5}{0.5/\sqrt{20}}
\]

Obliczamy liczbowo:
\[
	\sigma/\sqrt{20} = 0.5/\sqrt{20} \approx 0.5/4.472136 = 0.1118034
\]
\[
	a = \frac{-0.1}{0.1118034} = -0.8944272 \approx -0.8944
\]
\[
	b = 0
\]

Zatem:
\[
	P(8 \le Y \le 10) \approx P(-0.8944 \le Z \le 0),\quad Z\sim N(0,1)
\]
\[
	= \Phi(0) - \Phi(-0.8944) \approx 0.5 - 0.1856 = 0.3144
\]

\subsection*{3. CTG z poprawką ciągłości}

Poprawka: \( P(8 \le Y \le 10) = P(7.5 \le Y \le 10.5) \), czyli:
\[
	\left\{ \frac{7.5}{20} \le \bar{X}_{20} \le \frac{10.5}{20} \right\} = \{0.375 \le \bar{X}_{20} \le 0.525\}
\]

Podstawiamy do wzoru CTG:
\[
	a' = \frac{0.375 - 0.5}{0.5/\sqrt{20}} = \frac{-0.125}{0.1118034} = -1.118034 \approx -1.1180
\]
\[
	b' = \frac{0.525 - 0.5}{0.5/\sqrt{20}} = \frac{0.025}{0.1118034} = 0.2236068 \approx 0.2236
\]

Zatem:
\[
	P(8 \le Y \le 10) \approx P(-1.1180 \le Z \le 0.2236)
\]
\[
	= \Phi(0.2236) - \Phi(-1.1180) \approx 0.5885 - 0.1318 = 0.4567
\]
\subsection*{5. Porównanie}
Dokładna wartość: \( 0.4564 \).
CTG bez poprawki: \( 0.3144 \) (błąd \( \approx 0.1420 \)).
CTG z poprawką: \( 0.4567 \) (błąd \( \approx 0.0003 \)).

Przybliżenie CTG z poprawką ciągłości daje wynik bardzo bliski dokładnemu.
